% ==============================================================================
% SECTION: Recurrence of Failure and Repair Patterns
% ==============================================================================

\subsection{Recurrence of Failure and Repair Patterns}
\label{sec:benchmark-recurrence}

Historical repair experience can benefit future CI repair only when
related failure and repair patterns recur. We therefore examine whether
CI-Repair-Bench contains recurring repair instances and characterize
recurrence at two scopes: within individual repositories and across the
benchmark as a whole. To account for both explicit attribute overlap
and similarity in heterogeneous CI evidence, we measure recurrence from
two complementary views: \emph{structural similarity} and
\emph{lexical similarity}.

\paragraph{CI context normalization.}
Raw CI logs are heterogeneous and contain substantial execution noise.
We therefore first normalize the failure evidence of each instance into
a structured CI context using a fixed LLM-based extraction procedure.
Given the CI logs and workflow evidence, the extractor identifies
failure categories, error names, failure signals, validation commands,
tools, relevant packages or dependencies, and affected files. The same
prompt and output schema are applied to every benchmark instance.
Ground-truth repair information, including the developer patch and
modified files, is retained separately. This preprocessing provides a
consistent representation for comparing failure and repair
characteristics across heterogeneous CI executions.

\paragraph{Structural similarity.}
Structural similarity captures explicit overlap in the normalized CI
context and repair artifacts. For each pair of instances, we compute
Jaccard similarity over corresponding set-valued attributes, including
failure categories, error/failure signals, validation commands, tools,
packages or dependencies, and files associated with the failure and
developer repair. For attribute $k$, Jaccard similarity is

\begin{equation}
J_k(i,j)
=
\frac{|A_i^k \cap A_j^k|}
{|A_i^k \cup A_j^k|},
\end{equation}

where $A_i^k$ and $A_j^k$ denote the normalized values of attribute
$k$ for instances $i$ and $j$. We aggregate the available
attribute-level similarities to obtain the structural similarity
$S_{\mathrm{struct}}(i,j)$. This view captures recurring CI
characteristics and repair artifacts through explicit overlap.

\paragraph{Lexical similarity.}
Exact structural overlap may miss related instances expressed through
different error messages, workflow configurations, or patches. We
therefore complement structural similarity with lexical similarity.
For each instance, we construct a textual representation from its
normalized CI context, relevant failure-log evidence, workflow context,
and ground-truth repair information. We represent this text using
TF--IDF over unigrams and bigrams and compute cosine similarity,

\begin{equation}
S_{\mathrm{lex}}(i,j)
=
\frac{\mathbf{v}_i^\top \mathbf{v}_j}
{\|\mathbf{v}_i\|_2\|\mathbf{v}_j\|_2},
\end{equation}

where $\mathbf{v}_i$ denotes the TF--IDF representation of instance
$i$. This view captures recurring error signatures, commands,
dependency names, workflow terms, and repair-related tokens that may
not be reflected by exact attribute overlap.

\paragraph{Recurrence analysis.}
We quantify recurrence using nearest-neighbor similarity under two
candidate scopes. For an instance $i$, within-repository recurrence is
the similarity to its most similar distinct instance from the same
repository,

\begin{equation}
R_{\mathrm{within}}(i)
=
\max_{\substack{j \neq i \\ r_j = r_i}}
S(i,j),
\end{equation}

where $r_i$ denotes the repository containing instance $i$. Overall
benchmark recurrence is defined as the similarity to the most similar
distinct instance anywhere in CI-Repair-Bench,

\begin{equation}
R_{\mathrm{overall}}(i)
=
\max_{j \neq i} S(i,j).
\end{equation}

We compute both recurrence measures independently under structural
similarity, $S_{\mathrm{struct}}$, and lexical similarity,
$S_{\mathrm{lex}}$, and report their mean across eligible benchmark
instances. Within-repository recurrence characterizes repeated
project-specific failure--repair patterns, whereas overall recurrence
measures the extent to which related patterns occur anywhere in the
benchmark. Together, these analyses assess whether CI-Repair-Bench
contains recurring repair experience that can motivate historical
memory-based CI repair.

% ------------------------------------------------------------------------------

Table~\ref{tab:benchmark_recurrence} reports mean nearest-neighbor
similarity under both scopes. We observe moderate structural recurrence
within repositories (0.436), indicating that repositories exhibit
recurring combinations of failure categories, tools, and affected files.
Lexical recurrence is substantially higher (0.811 within-repository,
0.796 overall), suggesting that error messages, workflow configurations,
and repair patterns exhibit strong textual similarity even when explicit
structural overlap is partial. The similar magnitudes of within-repository
and overall benchmark recurrence (0.436 vs.\ 0.445 structural, 0.811 vs.\
0.796 lexical) indicate that recurring patterns occur both within individual
projects and across the broader benchmark. These findings validate the
premise that historical repair experience is available and exhibits
recurring patterns that can inform memory-based CI repair.

% ==============================================================================
% TABLE: Recurrence Similarity
% ==============================================================================

\begin{table}[t]
\caption{Recurrence of similar CI repair instances in CI-Repair-Bench.
         Values report mean nearest-neighbor similarity under two scopes:
         within-repository and overall benchmark.}
\label{tab:benchmark_recurrence}
\centering
\begin{tabular}{lcc}
\hline
\textbf{Scope} &
\textbf{Structural} &
\textbf{Lexical} \\
\hline
Within repository  & 0.436 & 0.811 \\
Overall benchmark  & 0.445 & 0.796 \\
\hline
\end{tabular}
\end{table}

% ==============================================================================
% IMPLEMENTATION NOTES (for appendix or supplementary material)
% ==============================================================================

\paragraph{Implementation details.}

\textbf{CI context extraction.}
The normalization uses GPT-4o (gpt-4o-2024-08-06) with temperature 0.0
and a fixed JSON schema. The extraction prompt instructs the model to
identify failure categories, error signals with tool names and locations,
validation commands, and affected files from the CI logs and workflow
evidence. All 543 extracted CI contexts are included in the benchmark
release. A manual validation on a stratified sample of 50 instances can
be conducted to assess extraction quality.

\textbf{File path handling.}
For within-repository comparisons, we use full file paths since path
structure is meaningful within a single codebase. For cross-repository
comparisons, we use basenames only, as full paths are not comparable
across different repositories.

\textbf{TF--IDF vectorization.}
We use scikit-learn's TfidfVectorizer with a vocabulary size of 1000,
English stop word removal, unigrams and bigrams (ngram\_range=(1,2)),
and minimum document frequency of 2.

\textbf{Affected vs.\ changed files.}
We distinguish between \emph{affected files} (mentioned in CI failure
evidence) and \emph{changed files} (modified by the developer repair
patch). The former describes failure localization; the latter describes
the actual repair. Both are included as separate attributes in
structural similarity computation.
